\section{The accepted-service expansion problem}\label{sec:r22-early}
The common object is the value of enlarging a contractually admissible decision class. For a context $\xi$, let $X$ be the accepted class, let $Y\subseteq X$ be a restricted class, and normalize the announced outside protocol to displacement zero. Write
\begin{equation}\label{eq:r22-common}
 J(\xi)=\max_{x\in X}F_\xi(x),\qquad
 K(\xi)=\max_{x\in Y}F_\xi(x),\qquad
 \Delta(\xi)=J(\xi)-K(\xi).
\end{equation}
Both classes contain zero. The institution determines these classes before optimization; they cannot be weakened when a learned proposal violates a commitment. The canonical sections below derive the outside and information-class comparators from service primitives. The continuation model then replaces root-only participation by every modeled continuation obligation. In its convex finite-tree implementation,
$X=\{x\in[\ell,\bar u]:Ax\leq a,\ Ex=0\}$; general vector and signed-coefficient formulations retain their actual rows rather than invoking positive-tariff transfer coordinates outside their scope.

The analysis answers three distinct questions about \eqref{eq:r22-common}. Continuation-price balance identifies when expansion is valuable. The resource-price response and its error bound identify which approximation errors matter for implementing that value. A final primal--dual audit determines whether the selected feasible decision meets the requested accuracy, while a restricted upper bound determines whether its gain exceeds an \emph{optimized} comparator. The last comparison is valid without assuming that the prediction method is the fastest solver. Our experiments report decision quality and complete computation separately, and retain both neural and direct-price implementations of the same accepted response.
